\documentclass[11pt]{article}
\usepackage{amsmath,amssymb}
\usepackage{geometry}
\geometry{a4paper,margin=1in}
\newcommand{\fn}[1]{{\small\texttt{\detokenize{#1}}}}
\newcommand{\eps}{\varepsilon}
\title{Attention Rationale Supervision: Training Objective}
\author{}
\date{}

\begin{document}
\maketitle

\section{Inputs and Supervision}
The following formulas describe a single function; sample indices are omitted.
\begin{itemize}
    \item \fn{data.py: prepare()} constructs vulnerable/patched function examples.
    Each retained function is encoded in full; overlength pairs are excluded.
    \item $x$: a function; $y=1$ for vulnerable and $y=0$ for patched.
    These are dataset labels, not guarantees that patched code is vulnerability-free.
    \item $V$: valid token-bearing source lines; $T_j$: tokens mapped to line $j$.
    Special tokens with line ID $-1$ do not contribute to line scores.
    \item $R\subseteq V$: changed/deleted lines on the vulnerable side,
    computed by\\ \fn{changed_vulnerable_lines()}.
    These are \textbf{proxy rationales}, not human-verified root-cause annotations.
    Patched examples have no positive rationale targets.
\end{itemize}

\section{Function Classification Loss}
\begin{itemize}
    \item \fn{model.py: Classifier.forward()} applies a two-class linear head to
    the final encoder layer's CLS hidden state $h^{(L)}_{\mathrm{CLS}}$:
    \[
        z=W_c h^{(L)}_{\mathrm{CLS}}+b_c,
        \qquad \pi=\operatorname{softmax}(z).
    \]
    \item Both vulnerable and patched functions contribute to cross-entropy:
    \[
        \boxed{\mathcal L_{\mathrm{cls}}(x,y)
        =-\log P_\theta(y\mid x)=-\log\pi_y.}
    \]
    \item Implementation: \fn{train.py: main()},
    \fn{torch.nn.functional.cross_entropy()}.
\end{itemize}

\section{CLS Attention to a Line Distribution}
\begin{itemize}
    \item \fn{model.py: Classifier.__init__()} selects layer $\ell$:
    the last global-attention layer for ModernBERT, or the last attention layer
    for RoBERTa/CodeBERT. It does not sum all layers.
    \item \fn{cls_probabilities()} computes the masked CLS-query attention
    probabilities and averages over $H$ heads:
    \[
        a_t=\frac{1}{H}\sum_{h=1}^{H}
        A^{(\ell,h)}_{\mathrm{CLS},t}.
    \]
    \item \fn{line_distribution()} averages tokens within each valid line,
    then normalizes over valid lines:
    \begin{align}
        u_j&=\frac{1}{|T_j|}\sum_{t\in T_j}a_t,\\
        p_j&=\frac{u_j}{\max(\sum_{k\in V}u_k,\eps)},
        \qquad \eps=10^{-12}.
    \end{align}
    \item Normally $\sum_{j\in V}p_j=1$; the denominator clamp protects
    numerical edge cases. These are relative line-attention scores,
    \textbf{not calibrated line-level vulnerability probabilities}.
\end{itemize}

\section{Rationale Annotation Loss}
\begin{itemize}
    \item The target is uniform across the annotated lines:
    \[
        q_j=\begin{cases}
            1/|R|,& j\in R,\\
            0,& j\in V\setminus R.
        \end{cases}
    \]
    \item Only vulnerable functions with nonempty valid targets participate
    ($y=1$ and $|R|>0$). For one such function,
    \[
        \boxed{\mathcal L_{\mathrm{rat}}(x,R)
        =-\sum_{j\in V}q_j\log p_j
        =-\frac{1}{|R|}\sum_{j\in R}\log p_j.}
    \]
    Implementation: \fn{model.py: rationale_loss()} clamps $p_j$ below
    at $\eps=10^{-12}$ before taking the logarithm.
    \item Each eligible function receives equal weight; within it, each
    annotated line receives equal weight. Patched examples have classification
    loss only, with no uniform-attention or zero-attention target.
    \item Ignoring numerical clamps, this is cross-entropy $H(q,p)$:
    \[
        H(q,p)=D_{\mathrm{KL}}(q\Vert p)+H(q),
        \qquad H(q)=\log|R|.
    \]
    Thus its gradients match forward KL, but the scalar loss includes the
    target-entropy constant. For two annotated lines, the per-function loss is
    $-\tfrac12(\log p_{j_1}+\log p_{j_2})$.
\end{itemize}

\section{Combined Loss and Fixed Scale Calibration}
\begin{itemize}
    \item The combined objective is
    \[
        \boxed{\mathcal L=
        \mathcal L_{\mathrm{cls}}
        +\lambda b\,\mathcal L_{\mathrm{rat}}.}
    \]
    \item In \fn{train.py: main()}, the classification term is averaged over
    all functions in the mini-batch; the rationale term is averaged separately
    over eligible vulnerable functions only. With no eligible functions, the
    rationale term is zero. The two averages can have different denominators;
    the combined objective above uses these averaged terms during training.
    \item $\lambda$ is \fn{--attention-weight}, default $0.1$.
    Gradients update the encoder and classification head; attention supervision
    is differentiable through the encoder's attention computations.
    \item Before optimization, sample at most 128 \emph{function examples}
    from train using the configured seed. In evaluation mode (dropout disabled),
    estimate the initial mean losses and set
    \[
        b=\frac{\overline{\mathcal L}^{(0)}_{\mathrm{cls}}}
        {\overline{\mathcal L}^{(0)}_{\mathrm{rat}}}.
    \]
    The bars denote arithmetic means of the initial calibration batch losses.
    Calibration batches without eligible rationales contribute zero to the
    denominator average. A nonpositive denominator raises an error.
    \item $b$ is fixed for all epochs and saved as \fn{loss_balance} in
    \fn{run.json}. At calibration, $\lambda b\,\overline{\mathcal L^{(0)}_{\mathrm{rat}}}
    =\lambda\overline{\mathcal L^{(0)}_{\mathrm{cls}}}$;
    this ratio is not held constant during training.
    \item Both losses are reduced within each mini-batch. Gradient accumulation
    averages these batch losses, rather than recomputing a single rationale
    mean over all eligible examples across the accumulated batches.
\end{itemize}

\section{Experimental Controls and Inference}
\begin{itemize}
    \item \fn{classification}: $\mathcal L=\mathcal L_{\mathrm{cls}}$ only.
    \item \fn{rationale}: the CLS-attention objective above.
    \item \fn{random}: same objective, but replace $R$ with fixed sampled
    lines of the same count, preferentially non-rationale lines. This is not
    a fully uniform random-line control.
    \item \fn{hidden}: replace the attention-derived distribution with
    \[
        s_j=w^\top\left(\frac{1}{|T_j|}
        \sum_{t\in T_j}h^{(L)}_{t}\right)+c,
        \qquad p_j=\operatorname{softmax}_{j\in V}(s_j).
    \]
    Use the same proxy targets and combined-loss structure, with its own
    calibration coefficient. Implemented by the linear token head and
    \fn{line_distribution(..., logits=True)}.
    \item Inference receives code only. Rationale labels are used for training,
    evaluation, and visualization, not as model inputs. Since line scores are
    normalized, even a patched function still has a highest-scoring line.
\end{itemize}

\end{document}
